\begin{enumerate}

  \item 
  
  等式的传递性表明，如果 $a=b$ 且 $b=c$，那么 $a=c$。蕴含箭头是否满足传递性？如果满足，请陈述之。
  \wbvfill
  
  \hint{
  
  
  我有时喜欢将蕴涵 $X \implies Y$ 改写为“X的真迫使Y为真。”这有帮助吗？如果我们知道X为真迫使Y为真，并且我们也知道Y为真将迫使Z为真，我们能得出什么结论？
  \vfill
  
  }
  
  \item 
  
  完成复合句 $A \implies B$ 和 ${\lnot}A \lor B$ 的真值表。
  \wbvfill
    
  \hint{
  
  你绝对应该能自己完成这个，但无论如何，这里是表格的框架：
  
  \begin{tabular}{|c|c|c|c|} \hline
  \rule[-6pt]{0pt}{24pt}  $A$ & $B$ & $A \implies B$ & ${\lnot}A \lor B$ \\ \hline
  \rule[-6pt]{0pt}{24pt}  $T$ &  $T$ & & \\ \hline
  \rule[-6pt]{0pt}{24pt}  $T$ & $\phi$ & & \\ \hline	 	 
  \rule[-6pt]{0pt}{24pt}  $\phi$ & $T$ & & \\ \hline
  \rule[-6pt]{0pt}{24pt}  $\phi$ & $\phi$ & & \\ \hline
  \end{tabular}
  
  \vfill
  
  }
  \workbookpagebreak
  
  \item 
  
  完成复合句 $A \implies (B \implies C)$ 和句子 $(A \implies B) \implies C$ 的真值表。关于条件句和结合律，你能得出什么结论？
  
  \wbvfill
  
  \hint{
  
  除了说明蕴涵{\bf 不}满足结合律之外，这个问题没有其他帮助。
  \vfill
  
  }
  
  
  \hintspagebreak
  
  \item 
  使用{\em 与}联结词($\land$)确定一个句子，该句子是 $A \implies B$ 的否定。
  \wbvfill
  
  \hint{
  
  嗯……这提示可能有点奇怪，但如果你在操场上听到一个孩子说：“哦是吗？我确实叫了你妈妈胖子，但你还没揍我！哎哟！嗷！！！别打了！！”
  
  他试图否定的是哪个条件句？
  }
  
  \item 
  
  将句子“修好马桶，否则我不付房租！”改写为条件句。
  \wbvfill
  
  \hint{
  在我看来，有八种可能的方式可以将“你修好马桶”和“我会付房租”（或它们各自的否定）围绕一个蕴涵箭头排列。这里是所有的方式。你来决定哪一个听起来最好。
  如果你修好马桶，那么我会付房租。\newline
  如果你修好马桶，那么我不会付房租。\newline
  如果你不修好马桶，我会付房租。\newline
  如果你不修好马桶，那么我不会付房租。\newline
  如果我付了房租，那么你一定修好了马桶。\newline
  如果我付了房租，那么你一定没有修好马桶。\newline
  如果我没有付房租，那么你一定修好了马桶。\newline
  如果我没有付房租，那么你一定没有修好马桶。\newline
  
  其中一些真的很奇怪……
  }
  
  \workbookpagebreak
  
  \item 
  为什么“如果猪会飞，我就是美索不达米亚的国王”这句话是真的？
  
  \wbvfill
  
  \hint{
  
  除非我们谈论的是某个名人把他们的宠物越南大肚猪带进头等舱，或者可能是某种弹射器……前件（如果部分）是假的，所以耶！我就是美索不达米亚的国王！！哇哦！什么？我不是？哦。但这个如果-那么句子是真的。真扫兴。
  }
  
  \item 
  使用皮尔斯箭头和/或谢弗竖线来表达陈述 $A \implies B$。（参见上一节的练习~\ref{ex:nand_nor}。）
  
  \wbvfill
  
  \hint{
  
  你会想使用 $\vert$，即谢弗竖线，也就是与非门，因为它的真值表包含三个T和一个$\phi$——你只需要弄清楚应该否定哪个输入，以使那个唯一的$\phi$出现在表的第二行而不是第一行。
  }
  
  
  
  \item 
  
  找出下列句子的逆否命题。
  \begin{enumerate}
    \item \wbitemsep 
    如果你承担不了后果，就不要做坏事。
    \item \wbitemsep 
    如果你在学校表现好，你就会找到一份好工作。
    \item \wbitemsep 
    如果你希望别人以某种方式对待你，你必须以那种方式对待别人。
    \item \wbitemsep 
    如果下雨，就一定有云。
    \item \wbitemsep 
  如果对于所有的 $n$ 都有 $a_n \leq b_n$，并且 $\sum_{n=0}^\infty b_n$ 是一个收敛级数，那么 $\sum_{n=0}^\infty a_n$ 也是一个收敛级数。
  \end{enumerate}
  
  %\wbvfill
  
  \hint{
  \begin{enumerate}
  \item 
  如果你做了坏事，你就必须承担后果。
  \item 
  如果你没有一份好工作，你一定是在学校表现不好。
  \item 
  如果你不以某种方式对待他人，你就不能指望他人以那种方式对待你。
  \item 
  如果没有云，就不可能下雨。
  \item 
  如果 $\sum_{n=0}^\infty a_n$ 不是一个收敛级数，那么要么对于某个 $n$ 有 $a_n > b_n$[注意：原文此处应为>], 要么 $\sum_{n=0}^\infty b_n$ 不是一个收敛级数。
  \end{enumerate}
  }
  \rule{0pt}{0pt}
  
  %\wbvfill
  
  \workbookpagebreak
  
  \item 
  “如果你掩护我，我也会掩护你”的逆命题和否命题是什么？
  \wbvfill
  
  \hint{
  
  逆命题是“如果我掩护你，那么你也会掩护我。”（听起来有点傻，不是吗——有点像一厢情愿……）
  否命题是“如果你不掩护我，那么我也不会掩护你。”（听起来不那么空洞，但意思是一样的……）
  }
  
  \item 
  
  微积分中的积分判别法用于判断一个无穷级数是收敛还是发散：假设 $f(x)$ 是一个正的、递减的实值函数，且 $\lim_{x \longrightarrow \infty} f(x) = 0$，如果反常积分 $\int_0^\infty f(x)$ 有一个有限值，那么无穷级数 $\sum_{n=1}^\infty f(n)$ 收敛。积分判别法可以通过将级数对应于积分的右黎曼和来形象地理解，因为函数是递减的，右黎曼和是积分值的下估计，因此
  
  \[ \sum_{n=1}^\infty f(n) < \int_0^\infty f(x).
  \]
  
  讨论积分判别法中条件陈述的否命题、逆命题和逆否命题的含义，并在可能的情况下为其提供理由。
  \wbvfill
  
  \hint{
 
  否命题说——如果积分不是有限的，那么级数就不收敛。你可以构造一个函数来证明这是错误的，例如，创建一个在整数x值之间有垂直渐近线的函数。即使只有一个这样的极点也足以使积分变为无穷大。
  逆命题说，如果级数收敛，那么积分必须是有限的。我们刚才讨论的反例在这里也适用。
  
  逆否命题说，如果级数不收敛，那么积分必定不是有限的。如果我们被允许使用不连续函数，不难想出一个其下方实际面积为零的 $f$ ——只需让f在除了整数x值之外的地方恒为零，在整数x值处取与级数项相同的值即可。但是等等，我们刚才描述的函数不是“递减”的——这可能就是为什么把那个假设放在那里的原因！
  }
  
  \rule{0pt}{0pt}
  
  \wbvfill
  
  \workbookpagebreak
  
  \item 
  在骑士与无赖岛上（见第~\pageref{IKK}页），你遇到了两个名叫洛克和德摩斯梯尼的人。
  洛克说：“德摩斯梯尼是个无赖。”\newline
  德摩斯梯尼说：“洛克和我是骑士。”
  
  谁是骑士，谁是无赖？
  \wbvfill
  
  \hint{
  德摩斯梯尼可能在说真话吗？}
  
  \end{enumerate}